% Unite: papier 34, titre repete, chapitre III et section 15; traduction francaise depuis la source allemande auditee.

\providecommand{\mo}{\mathfrak{o}}
\providecommand{\mM}{\mathfrak{M}}

\subsection*{Chapitre III. Théorie des modules et des représentations}

\subsection*{\S{} 15. Représentations et modules de représentation\textsuperscript{15)}}
Soit $\mo$ un anneau, et $K$ un anneau avec élément unité. (Dans les applications ultérieures, $K$ est toujours un corps.)

Une représentation de degré $n$ de $\mo$ dans $K$ est une homomorphie
\[
        \mo\sim\mathfrak D,
\]
où $\mathfrak D$ est un anneau de matrices de degré $n$ sur $K$.

Par module de représentation de $\mo$ relativement à $K$, on entend un bimodule $\mM$ (\S{}1, 5), qui est module à gauche sur $\mo$ et module à droite sur $K$:
\[
        \mo\mM\subseteq\mM,
        \qquad
        \mM K\subseteq\mM,
\]
qui est en outre somme directe d'un nombre fini de modules $K$ monogènes:
\[
        \mM=x_1K+\cdots+x_nK;
\]
et où l'élément unité de $K$ est opérateur unité.

Tout module de représentation conduit à une représentation. Soit $c$ dans $\mo$ et
\[
        cx_k=\sum_i x_i\gamma_{ik},
\]
ou encore
\[
        c(x_1,\ldots,x_n)=(x_1,\ldots,x_n)C,
        \qquad C=(\gamma_{ik}).
\]
Alors les matrices $C$ forment une représentation pour $c$, car
\[
        (b+c)x_k=\sum_i x_i(\beta_{ik}+\gamma_{ik}),
\]
\[
\begin{aligned}
        bcx_k
        &=b\sum_jx_j\gamma_{jk}
          =\sum_j b x_j\gamma_{jk}
          =\sum_j\sum_i x_i\beta_{ij}\gamma_{jk}  \\
        &=\sum_i x_i\biggl(\sum_j\beta_{ij}\gamma_{jk}\biggr),
\end{aligned}
\]
ou
\[
        bc(x_1,\ldots,x_n)=(x_1,\ldots,x_n)BC.
\]

Réciproquement: toute représentation $\mathfrak D$ de $\mo$ appartient à un module de représentation, et même à une base déterminée de celui-ci. On entend en effet par $\mM$ l'ensemble des formes linéaires formelles en $x_1,\ldots,x_n$,
\[
        y=\sum_i x_i\alpha_i.
\]
Alors $\mM$ est un module à droite sur $K$. Définissons en outre, lorsque la matrice $C=(\gamma_{ik})$ est associée à l'élément $c$,
\[
        cx_k=\sum_i x_i\gamma_{ik},
        \qquad
        c\biggl(\sum_k x_k\alpha_k\biggr)=\sum_i x_i\sum_k\gamma_{ik}\alpha_k.
\]
Des relations d'homomorphie
\[
        c+d\to C+D,
        \qquad cd\to CD
\]
on obtient
\[
        (c+d)x_i=cx_i+dx_i,
        \qquad cdx_i=c(dx_i),
\]
et il en va de même pour les sommes $\sum_i x_i\alpha_i$. Les autres propriétés de bimodule
\[
        c(y+z)=cy+cz,
        \qquad c(y\alpha)=(cy)\alpha
\]
sont triviales. Ainsi $\mM$ est un module de représentation qui, par construction, appartient exactement à la représentation $\mo\sim\mathfrak D$.

Soit maintenant $(y_1,\ldots,y_n)$ une autre base du même module de représentation:
\[
        (y_1,\ldots,y_n)=(x_1,\ldots,x_n)P,
        \qquad
        (x_1,\ldots,x_n)=(y_1,\ldots,y_n)P^{-1}.
\]
Si $b\sim B$ relativement à la base $x$, on a
\[
        b(y_1,\ldots,y_n)=b(x_1,\ldots,x_n)P
        =(x_1,\ldots,x_n)BP
        =(y_1,\ldots,y_n)P^{-1}BP.
\]
On appelle les représentations $b\sim B$ et $b\sim P^{-1}BP$ équivalentes, et on les compte dans la même classe de représentations. Comme toute matrice régulière $P$ transforme de nouveau une base de $\mM$ en une base, il est démontré:

\emph{Tout module de représentation conduit de manière univoque à une classe de représentations déterminée.}\textsuperscript{15a)}

{\footnotesize\noindent 15) Les modules de représentation relativement à des corps commutatifs se trouvent d'abord chez W. Krull, \emph{Theorie und Anwendung der verallgemeinerten Abelschen Gruppen}, Sitzungsber. Heidelberger Ak. 1926, 1. Abhandl.\par
15a) Addition lors de la correction (14 avril 1929). Comme B. L. van der Waerden me le communique, on peut obtenir directement un lien invariant, indépendant du choix spécial de la base, en séparant les notions: transformation linéaire et matrice. Une transformation linéaire est un homomorphisme de deux modules de formes linéaires; une matrice est l'expression (la représentation) de cet homomorphisme pour un choix déterminé de base. I. Tout module de représentation associe univoquement au domaine des multiplicateurs $\mo$ un système homomorphe de transformations linéaires du module en lui-même. En effet, d'après la fin du \S{}2, les multiplicateurs à gauche d'un bimodule $\mM$ engendrent des homomorphismes d'opérateurs de $\mM$ en lui-même relativement aux opérateurs à droite (ici la multiplication par $K$); et la correspondance est homomorphe. II. Tout système, homomorphe à $\mo$, de transformations linéaires d'un module de formes linéaires en lui-même conduit à un module de représentation. En effet, d'après encore la fin du \S{}2, ces homomorphismes, pris comme multiplicateurs à gauche, donnent un bimodule; il en va donc de même relativement au domaine des multiplicateurs $\mo$ si l'on pose $cm=\Gamma m$ pour tout $m$ de $\mM$ et tout $c$ de $\mo$, où $c\to\Gamma$ est l'homomorphisme d'anneaux initialement donné. Aux modules de représentation isomorphes comme bimodules -- c'est-à-dire par un isomorphisme compatible avec les actions à droite et à gauche -- correspondent les mêmes transformations, et aux modules qui ne sont pas isomorphes comme bimodules des transformations différentes. Si, pour un choix déterminé de base, on exprime chaque transformation par une matrice, I donne la représentation $\mo\sim\mathfrak D$, tandis que II affirme que toute représentation de cette sorte appartient à un module de représentation et à une base déterminée de celui-ci. On retrouve ainsi, sans aucun calcul, la correspondance entre modules de représentation et classes de représentations.\par}

Il est clair que deux modules de représentation isomorphes au sens des opérateurs correspondent à la même classe de représentations. Mais réciproquement aussi: si deux modules de représentation $(x_1,\ldots,x_n)$ et $(\bar x_1,\ldots,\bar x_n)$ engendrent la même représentation, alors la correspondance
\[
        x_i\to\bar x_i,
        \qquad
        \sum_i x_i\alpha_i\to\sum_i\bar x_i\alpha_i
\]
fournit un isomorphisme. Car la règle de multiplication est entièrement déterminée par les matrices $C$.

Cas particulier: si deux bases de $\mM$ engendrent la même représentation, elles se déduisent l'une de l'autre par un isomorphisme d'opérateurs de $\mM$ sur lui-même.

Ou encore: si $C=PCP^{-1}$ pour tous les $C$ et pour une matrice $P$ fixe, alors la correspondance
\[
        y_i\to x_i,
        \qquad y_i=\sum_kx_k\pi_{ik},
\]
est un isomorphisme de $\mM$ sur lui-même. Réciproquement: tout isomorphisme $y_i\to x_i$ du bimodule $\mM$ sur lui-même est réalisé par une matrice $P$ qui commute avec tous les $C$.

On peut aussi étendre la notion de représentation aux anneaux sur lesquels un domaine commutatif de multiplicateurs $P$ agit centralement (\S{}8). Dans ce cas on ne prend que des anneaux de représentation $K$ qui contiennent $P$ dans leur centre, et l'on exige de la représentation non seulement qu'elle soit un homomorphisme d'anneaux $\mo\sim\mathfrak D$, mais même qu'elle soit un homomorphisme d'opérateurs: de $a\to A$ doit suivre
\[
        a\rho\to A\rho\qquad(\rho\in P).
\]
Pour les modules de représentation, cette exigence signifie la règle de calcul supplémentaire
\[
        am\cdot\rho=a(m\rho)=a\rho\cdot m.
\]
L'action de $P$ sur le module commute alors avec celle de l'anneau.
